\clearpage
\section{Discussion}

The decomposition of $g(\tau)$'s variability into within- and between-individual components --- captured here by the smoothness parameter $\psi$ --- has a substantial impact on both uncontrolled and controlled epidemic dynamics. Narrow individual infectiousness profiles ($\psi \rightarrow 0$) inflate early-epidemic stochasticity, reduce the effective information content of contact tracing data, introduce overdispersion in secondary infection counts when contact rates vary over time, and make detect-and-isolate interventions sensitive to small shifts in detection timing. These effects are invisible to deterministic, mean-field models: the basic reproduction number $R_0$, the generation interval distribution $g(\tau)$, and the population-level infectiousness profile $A(\tau)$ are all unchanged by $\psi$. The effects emerge only in the stochastic dynamics, and can be substantial in magnitude. To study them, we introduced the Gamma time-shifted model, a flexible decomposition of $g(\tau)$ into shared and individual components, and combined it with a general contact process framework that extends the renewal equation to the individual level. Together, these tools provide a parsimonious way to assess an underexplored axis of uncertainty in epidemic modeling --- one that we expect will have implications for contact tracing, projection models, and intervention planning, even if only through sensitivity analyses.

A natural question is why the shape of the individual infectiousness profile has received so little attention. One reason is that $\psi$ does not affect mean-field dynamics, by design: $E_i[g_i(\tau)] = g(\tau)$ regardless of $\psi$, so standard renewal equation models, which depend only on $R_0$ and $g(\tau)$, are $\psi$-invariant. Consequently, epidemic forecasts and retrospective analyses that rely on the renewal equation can proceed without specifying $\psi$, and for many purposes this is adequate. A second reason is that the individual infectiousness profile has been dealt with implicitly in many settings, without being isolated as an independent axis of variation. Individual-based models, for example, routinely make assumptions about $a_i(\tau)$'s shape --- whether through SEIR-type stepwise profiles, viral kinetics-informed curves, or agent-based transmission rules --- but these assumptions are typically adopted for modeling convenience rather than studied as objects of interest. Similarly, the linear chain trick can be used to narrow the individual infectiousness profile by stringing together sequences of exposed and infectious compartments, but this simultaneously changes $g(\tau)$, conflating two distinct effects \citep{wearing_appropriate_2005, lloyd_realistic_2001}. The central contribution of the present work is to decouple these effects: the Gamma time-shifted model varies the width of $a_i(\tau)$ while holding $g(\tau)$ constant, enabling a clean assessment of the individual infectiousness profile's independent impact.

Several of our findings merit further discussion. For the early-epidemic time shift (Section~\ref{sec:results_time-shift}), the shift in $E[\varsigma]$ is modest; the main effect is the inflation of $\text{Var}[\varsigma]$, which widens the distribution of possible epidemic trajectories. This variability is relevant for scenario planning: early in an emerging epidemic, punctuated profiles make it harder to distinguish a ``slow'' epidemic from a ``late'' one, complicating early assessments of epidemic severity and the decision of when to intervene.

For inference of $r$ and $g(\tau)$ (Section~\ref{sec:results_r-g-uncertainty}), the increased overdispersion in daily case counts and the reduced effective sample size of serial interval data both point to the need for statistical methods that account for within-cluster correlation. Treating serial intervals as independent draws from $g(\tau)$ --- as is common practice --- can produce overconfident estimates of $\alpha$ and $\beta$ when individual infectiousness profiles are narrow. This suggests that contact tracing protocols should, where possible, record the identity of the index case for each observed serial interval, enabling cluster-level analyses that correctly account for the correlation structure.

The interaction between $\psi$ and time-varying contacts (Section~\ref{sec:results_overdispersion}) reveals a fundamentally distinct mechanism of superspreading. Classical overdispersion in secondary infection counts arises from ``super-shedders'' (variation in biological infectiousness $b_i$) or ``super-contactors'' (variation in contact levels $E_t[c_i(t)]$), both of which are in principle identifiable as individual-level traits. The overdispersion generated by the $g_i$--$c_i$ interaction, by contrast, arises from the coincidence of a person's narrow infectiousness window with a period of high or low contact intensity --- a form of ``bad luck'' that is not attributable to any persistent individual characteristic. This suggests a limit to which superspreading can be targeted by identifying high-risk individuals: even in a population of identical shedders with identical average contact rates, narrow infectiousness profiles can generate substantial overdispersion if contacts vary over time. Quantifying the relative contributions of these mechanisms to observed overdispersion is an important avenue for future work.

For detect-and-isolate interventions (Section~\ref{sec:results_detect-and-isolate}), the sensitivity of testing effectiveness to $\psi$ has practical implications. Shifts in symptom timing relative to peak infectiousness are commonplace --- as occurred across SARS-CoV-2 variants --- and changes in testing frequency are routine operational decisions. Our results indicate that the epidemic consequences of such shifts depend on the shape of the individual infectiousness profile: for punctuated profiles, small changes in detection timing can produce sharp, nonlinear changes in testing effectiveness, generation interval truncation, and overdispersion. This sensitivity underscores the importance of characterizing $a_i(\tau)$ for operational planning. Additionally, the differential response of testing effectiveness to detection timing across $\psi$-values suggests that intervention data could, in principle, provide information about the individual infectiousness profile beyond what is available from contact tracing alone --- an idea we have not developed here but consider a promising direction.

More broadly, the effects of detect-and-isolate interventions and gathering size restrictions on epidemic dynamics extend beyond their impact on $R_\text{eff}$. Detect-and-isolate interventions simultaneously reduce $R_\text{eff}$, truncate the generation interval (accelerating transmission among those who escape detection), and modulate overdispersion --- and these effects can work in opposing directions. Gathering size restrictions reduce $R_\text{eff}$ but also reduce overdispersion, which can lower epidemic extinction probabilities by a different amount than what the $R_\text{eff}$ reduction alone would imply. Evaluating interventions purely in terms of their effect on $R_\text{eff}$ can therefore be misleading; the framework introduced here provides a more complete accounting of how interventions reshape the full distribution of secondary infections, not just its mean.

Our simulation-based power analysis suggests that $\psi$ is identifiable from contact tracing data under reasonably favorable conditions, particularly for pathogens with moderate-to-high $R_0$. Our empirical estimates, drawn from published transmission trees for eight pathogens, suggest a spectrum of $\psi$-values: measles and pneumonic plague appear to have strongly punctuated profiles ($\hat{\psi} \approx 0$), while COVID-19, MERS, and smallpox appear more diffuse ($\hat{\psi} > 0.7$). There are biological reasons to expect such differences: some pathogens produce brief, intense bursts of shedding, while others sustain infectiousness over longer periods, and our estimates broadly align with this taxonomy.

However, these empirical estimates should be interpreted with considerable caution. Several biases in the available data are likely to push $\hat{\psi}$ downward. Contact tracing studies preferentially capture large clusters and superspreading events, which can arise from event-based transmission (\eg, a gathering at which many people are simultaneously exposed) rather than from a genuinely narrow infectiousness window; this ascertainment bias is particularly concerning for measles and pneumonic plague, where transmission is often event-mediated and well-traced. Incomplete observation of an infector's full set of secondary infections (due to imperfect ascertainment or right-censoring in ongoing outbreaks) can reduce within-cluster variance, mimicking the effect of low $\psi$. Date discretization to the level of symptom onset day compresses within-cluster variance further. Finally, publication bias in the OutbreakTrees database \citep{taube_outbreaktrees_2022} may over-represent notable outbreaks that tend to feature unusual clustering. Conversely, some biases push $\hat{\psi}$ upward: uncertainty in transmission direction can corrupt cluster assignments, and observation noise from variable incubation periods inflates apparent within-cluster variance. Additionally, our inference conditions on published generation interval parameters, introducing sensitivity to those estimates. We also note that local susceptible depletion could produce patterns that mimic low $\psi$: for a highly infectious pathogen like measles, an index case may rapidly exhaust local susceptibles upon reaching peak infectiousness, yielding tightly clustered secondary infections even if the biological window of infectiousness is broad. Disentangling such epidemiological effects from the intrinsic shape of $a_i(\tau)$ would require data from settings where susceptible depletion is minimal, such as controlled exposure studies or outbreaks in large, well-mixed populations.

Several additional limitations should be noted. The Gamma time-shifted model, while flexible, is only one possible parameterization of the within- \vs between-individual decomposition. Other families of distributions, or nonparametric approaches, might capture features that the Gamma model cannot. The contact process models we employ --- periodic and stochastic --- are useful for developing intuition but are difficult to calibrate empirically; real contact patterns are structured in ways (\eg, household \vs community \vs workplace contacts) that these simple models do not capture. Our treatment of detect-and-isolate interventions, while covering several detection mechanisms, does not address the full range of possibilities (\eg, contact tracing-initiated testing, adaptive screening). And our empirical inference of $\psi$ is intentionally simplified: a more rigorous approach would jointly estimate $\psi$ and the generation interval parameters, account for observation noise and ascertainment explicitly, and propagate uncertainty in incubation period estimates.

Despite these limitations, we believe this work opens productive avenues for research. The individual infectiousness profile is a fundamental epidemiological quantity that has been largely overlooked as an independent axis of variation, and our results demonstrate that it can meaningfully impact epidemic dynamics and intervention effectiveness. Incorporating $\psi$ --- even as a sensitivity parameter --- into projection and forecasting models would make those models more honest about a source of uncertainty that is currently ignored, and clearer about a mechanism that could drive some of the variability observed in real epidemics. Better data are needed: longitudinal studies of within-host viral kinetics linked to transmission outcomes, household studies that capture the full set of an infector's secondary cases, and controlled exposure experiments could all inform estimates of $a_i(\tau)$'s shape. We hope that the framework introduced here will motivate the collection of such data and support more nuanced assessments of epidemic risk and intervention impact.

